\begin{definition}
    Метрическое пространство $X$ называется \textit{компактным},
    если для любой системы открытых множеств
    $\{G_\alpha\}_{\alpha \in A} \subset 2^X$ такой,
    что $\bigcup\limits_{\alpha \in A} G_\alpha = X$,
    существует конечный набор $\alpha_1, \dotsc, \alpha_n \in A$ такой,
    что $\bigcup\limits_{k = 1}^n G_{\alpha_k} = X$.
\end{definition}

\begin{note}
    Определение выше можно сформулировать так: из любого открытого покрытия пространства $X$ можно выделить конечное подпокрытие. Или, эквивалентно, если система открытых множеств не содержит конечного покрытия пространства $X$, то она не является покрытием пространства $X$.
\end{note}

\begin{example}
    Пространство из замкнутых ограниченых множеств $\R^n$ будет компактным
\end{example}

\begin{definition}\label{center-system}
    Пусть $X$ "--- метрическое пространство. Система $\{B_\alpha\}_{\alpha \in A} \subset 2^X$ называется центрированной, если для любого конечного набора $\alpha_1, \dotsc, \alpha_n \in A$ выполнено $\bigcap\limits_{k = 1}^n B_{\alpha_k} \neq \varnothing$.
\end{definition}

\begin{theorem}\label{thm3.1}
    Метрическое пространство $X$ компактно $\lra$ любая центрированная система замкнутых множеств в $X$ имеет непустое пересечение.
\end{theorem}

\begin{proof}
    Каждой системе открытых множеств $\{G_\alpha\}_{\alpha \in A} \subset 2^X$ можно поставить в соответствие систему замкнутых множеств $\{F_\alpha\}_{\alpha \in A} := \{X \bs G_\alpha\}_{\alpha \in A}$, и наоборот. Тогда любая система открытых множеств $\{G_\alpha\}_{\alpha \in A}$, не содержащая конечного покрытия, не является покрытием $\lra$ любая центрированная система замкнутых множеств $\{F_{\alpha}\}_{\alpha \in A}$ имеет непустое пересеченеие.
\end{proof}